\usepackage{graphicx}
\usepackage{mathtools}
\usepackage{amssymb}
\usepackage{amsthm}
\usepackage{thmtools}
\usepackage{tocloft}
\usepackage{listings}
\usepackage{xcolor}
\usepackage{eso-pic}
\usepackage{transparent}
\usepackage{fontawesome5}
\usepackage{nameref}
\usepackage{babel}
\usepackage{hyperref}
\usepackage{fancyhdr}
\usepackage{enumitem}
\usepackage{caption}
\usepackage{float}
\usepackage{multirow}
\usepackage{pmboxdraw}

% DEFINIR VARIABLES DEL TEMA (Antes de usarlas)
\newcommand{\asignaturename}{Programación para Ciencia de Datos} % <-- Cambia esto para el nombre de la asignatura
\newcommand{\coursenumber}{2º } % <-- Cambia esto para el curso
\newcommand{\courseyear}{2025/2026} % <-- Cambia esto para el año académico

% Definimos la marca de agua
\newcommand\BackgroundPic{%
  \put(0,0){%
    \parbox[b][\paperheight]{\paperwidth}{
      \vfill
      \centering
      {\transparent{0.05}\rotatebox{-45}{\fontsize{80}{100}\selectfont\textcolor{gray}{\textbf{Alfonso Andrés}}}}%
      \vfill
    }
  }
}

% Aplicamos la marca de agua a todas las páginas excepto la portada
\pagestyle{fancy} 
\fancyhf{}
\fancyhead[L]{\textbf{\asignaturename}}
\fancyhead[R]{\textbf{Entregable I}}
\fancyfoot[L]{\textbf{Grado en Ciencia e Ingeniería de Datos (UPCT) - \coursenumber Curso - \courseyear}}
\fancyfoot[R]{\textbf{\thepage}}

% Redefinimos las líneas de encabezado y pie de página
\renewcommand{\headrulewidth}{0.4pt}
\renewcommand{\footrulewidth}{0.4pt}

% Aplicamos el mismo estilo a las páginas especiales
\fancypagestyle{plain}{
  \fancyhf{}
  \fancyhead[L]{\textbf{\asignaturename}}
  \fancyhead[R]{\textbf{Entregable I}}
  \fancyfoot[L]{\textbf{Grado en Ciencia e Ingeniería de Datos (UPCT) - \coursenumber Curso - \courseyear}}
  \fancyfoot[R]{\textbf{\thepage}}
  \renewcommand{\headrulewidth}{0.4pt}
  \renewcommand{\footrulewidth}{0.4pt}
}

% Guardamos el comando original de maketitle antes de redefinirlo
\let\oldmaketitle\maketitle

% Redefinimos el comando maketitle para tener una portada limpia más una segunda página en blanco más restaurar estilos después
\renewcommand{\maketitle}{%
  \ClearShipoutPictureBG% 1. Asegurar que no hay marca de agua
  \oldmaketitle% 2. Generar el título (Página 1)
  \thispagestyle{empty}% 3. Quitar encabezados/pies de la portada
  \clearpage% 4. Salto de página (Fin de la portada)

  % --- INICIO DE PÁGINA 2 EN BLANCO ---
  \thispagestyle{empty}% 5. Quitar encabezados/pies de la segunda página
  \null% 6. Poner algo invisible (necesario para que LaTeX cree la página)
  \clearpage% 7. Salto de página (Fin de la página en blanco)
  % --- FIN DE PÁGINA 2 EN BLANCO ---

  \AddToShipoutPictureBG{\BackgroundPic}% 8. Restaurar marca de agua (Página 3 en adelante)
  % Nota: Si quieres que la numeración empiece en 1 en la pág 3, descomenta la siguiente línea:
  % \setcounter{page}{1} 
}

% Guardamos el comando original de tableofcontents antes de redefinirlo
\let\oldtableofcontents\tableofcontents

% Redefinimos el comando tableofcontents para añadir un salto de página al final
\renewcommand{\tableofcontents}{\oldtableofcontents\clearpage}

% Editamos la profundidad de los apartados en la tabla de contenidos
\setcounter{tocdepth}{5}

% Añadir puntos suspensivos (leaders) a todos los niveles en la tabla de contenidos
\renewcommand{\cftsecleader}{\cftdotfill{\cftdotsep}}
\renewcommand{\cftsubsecleader}{\cftdotfill{\cftdotsep}}
\renewcommand{\cftsubsubsecleader}{\cftdotfill{\cftdotsep}}
\renewcommand{\cftparaleader}{\cftdotfill{\cftdotsep}}
\renewcommand{\cftsubparaleader}{\cftdotfill{\cftdotsep}}

% Ajustar el espacio entre el número y el título en todos los niveles
\setlength{\cftsecnumwidth}{3em}
\setlength{\cftsubsecnumwidth}{4em}
\setlength{\cftsubsubsecnumwidth}{5em}
\setlength{\cftparanumwidth}{6em}
\setlength{\cftsubparanumwidth}{7em}

% Poner en negrita TODAS las entradas del índice en todos los niveles
\renewcommand{\cftsecfont}{\bfseries}
\renewcommand{\cftsecpagefont}{\bfseries}
\renewcommand{\cftsubsecfont}{\bfseries}
\renewcommand{\cftsubsecpagefont}{\bfseries}
\renewcommand{\cftsubsubsecfont}{\bfseries}
\renewcommand{\cftsubsubsecpagefont}{\bfseries}
\renewcommand{\cftparafont}{\bfseries}
\renewcommand{\cftparapagefont}{\bfseries}
\renewcommand{\cftsubparafont}{\bfseries}
\renewcommand{\cftsubparapagefont}{\bfseries}

% Definimos el estilo visual de los teoremas
\newtheoremstyle{estilosinnumero}
  {\topsep}                       % Espacio vertical ANTES del entorno
  {\topsep}                       % Espacio vertical DESPUÉS del entorno
  {\normalfont\leftskip=2em}      % FUENTE CUERPO: Sangría de 2em para todo el texto
  {-2em}                          % SANGRÍA TÍTULO: Restamos 2em para pegar el título al margen
  {\bfseries}                     % FUENTE TÍTULO: Negrita
  {}                              % PUNTUACIÓN: Vacío
  {0pt}                           % SEPARACIÓN: 0pt (Necesario para evitar errores de compilación)
  {%
    \thmname{#1}\thmnote{ (#3)}   % Imprimir "Teorema (Nombre)"
    \\[0.5em]                     % Salto de línea + espacio vertical de 0.5em antes del cuerpo
  }

% Aplicamos el estilo
\theoremstyle{estilosinnumero}

% Creamos los entornos no enumerados
\newtheorem*{teorema}{Teorema}
\newtheorem*{proposicion}{Proposición}
\newtheorem*{lema}{Lema}
\newtheorem*{corolario}{Corolario}
\newtheorem*{definicion}{Definición}
\newtheorem*{ejemplo}{Ejemplo}
\newtheorem*{observacion}{Observación}
\newtheorem*{nota}{Nota}
\newtheorem*{demostracion}{Demostración}

% --- CONFIGURACIÓN DE CAPTIONS (TÍTULOS) Y CÓDIGO ---

% 1. Traducción de etiquetas al español
\renewcommand{\lstlistingname}{Code}

% 2. Estilo UNIFICADO para Tablas y Códigos
% Esto hace que ambos se vean idénticos (Centrados, Etiqueta en Negrita)
\captionsetup{
    labelfont=normalfont,        % "Tabla 1" o "Código 1" en Negrita
    textfont=normalfont, % El texto de la descripción en fuente normal
    format=plain,        % Formato estándar
    justification=centering, % Centrado (igual que en tu captura)
    labelsep=colon,      % Separador de dos puntos ": "
    figurename=Figura,   % Por si usas imágenes
    tablename=Tabla
}

% 3. Definición de colores para el código
\definecolor{backcolour}{rgb}{0.96, 0.96, 0.96}
\definecolor{codegreen}{rgb}{0, 0.6, 0}
\definecolor{codegray}{rgb}{0.5, 0.5, 0.5}
\definecolor{codepurple}{rgb}{0.58, 0, 0.82}
\definecolor{codeblue}{rgb}{0, 0, 0.8}

% 4. Configuración del entorno de código (Listings)
\lstdefinestyle{academico}{
    backgroundcolor=\color{backcolour},
    commentstyle=\color{codegreen},
    keywordstyle=\color{codeblue}\bfseries,
    numberstyle=\tiny\color{codegray},
    stringstyle=\color{codepurple},
    basicstyle=\ttfamily\fontsize{7}{9}\selectfont,
    breakatwhitespace=false,
    breaklines=true,
    % --- AQUÍ ESTÁ EL TRUCO PARA QUE PAREZCA UNA TABLA ---
    captionpos=b,        % 'b' = bottom (Abajo). Pon 't' si prefieres Arriba.
    abovecaptionskip=10pt, % Espacio entre el código y el título
    % -----------------------------------------------------
    keepspaces=true,
    numbers=none,
    numbersep=5pt,
    showspaces=false,
    showstringspaces=false,
    showtabs=false,
    tabsize=2,
    frame=single,
    rulecolor=\color{black!50},
    literate={ñ}{{\~n}}1 {á}{{\'a}}1 {é}{{\'e}}1 {í}{{\'i}}1 {ó}{{\'o}}1 {ú}{{\'u}}1
}

\lstset{style=academico}
% --- FIN CONFIGURACIÓN ---